
\documentclass{article}
\usepackage{colortbl}
\usepackage{makecell}
\usepackage{multirow}
\usepackage{supertabular}

\begin{document}

\newcounter{utterance}

\centering \large Interaction Transcript for game `hot\_air\_balloon', experiment `air\_balloon\_survival\_en\_complexity\_hard', episode 3 with Qwen3{-}32B{-}t0.0.
\vspace{24pt}

{ \footnotesize  \setcounter{utterance}{1}
\setlength{\tabcolsep}{0pt}
\begin{supertabular}{c@{$\;$}|p{.15\linewidth}@{}p{.15\linewidth}p{.15\linewidth}p{.15\linewidth}p{.15\linewidth}p{.15\linewidth}}
    \# & $\;$A & \multicolumn{4}{c}{Game Master} & $\;\:$B\\
    \hline

    \theutterance \stepcounter{utterance}  
    & & \multicolumn{4}{p{0.6\linewidth}}{
        \cellcolor[rgb]{0.9,0.9,0.9}{
            \makecell[{{p{\linewidth}}}]{
                \texttt{\tiny{[P1$\langle$GM]}}
                \texttt{You are participating in a collaborative negotiation game.} \\
\\ 
\texttt{Together with another participant, you must agree on a single set of items that will be kept. Each of you has your own view of how much each item matters to you (importance). You do not know how the other participant values the items. Additionally, you are given the effort each item demands.} \\
\texttt{You may only agree on a set if the total effort of the selected items does not exceed a shared limit:} \\
\\ 
\texttt{LIMIT = 6938} \\
\\ 
\texttt{Here are the individual item effort values:} \\
\\ 
\texttt{Item effort = \{"A08": 997, "B95": 950, "A39": 272, "A65": 448, "B81": 382, "B34": 804, "A28": 237, "C82": 56, "B69": 105, "A87": 612, "B64": 528, "B15": 527, "C79": 524, "B16": 167, "C53": 133, "C83": 300, "C10": 945, "A97": 50, "A25": 906, "C84": 70, "A71": 224, "B23": 4, "B96": 689, "C51": 64, "C17": 434, "B19": 750, "C72": 733, "A94": 873, "A46": 22, "B43": 68, "B37": 57, "A88": 10, "A96": 36, "C60": 551, "A75": 348\}} \\
\\ 
\texttt{Here is your personal view on the importance of each item:} \\
\\ 
\texttt{Item importance values = \{"A08": 1096, "B95": 1049, "A39": 371, "A65": 547, "B81": 481, "B34": 903, "A28": 336, "C82": 155, "B69": 204, "A87": 711, "B64": 627, "B15": 626, "C79": 623, "B16": 266, "C53": 232, "C83": 399, "C10": 1044, "A97": 149, "A25": 1005, "C84": 169, "A71": 323, "B23": 103, "B96": 788, "C51": 163, "C17": 533, "B19": 849, "C72": 832, "A94": 972, "A46": 121, "B43": 167, "B37": 156, "A88": 109, "A96": 135, "C60": 650, "A75": 447\}} \\
\\ 
\texttt{Goal:} \\
\\ 
\texttt{Your goal is to negotiate a shared set of items that benefits you as much as possible (i.e., maximizes total importance to YOU), while staying within the LIMIT. You are not required to make a PROPOSAL in every message {-} you can simply negotiate as well. All tactics are allowed!} \\
\\ 
\texttt{Interaction Protocol:} \\
\\ 
\texttt{You may only use the following structured formats in a message:} \\
\\ 
\texttt{PROPOSAL: \{'A', 'B', 'C', …\}} \\
\texttt{Propose keeping exactly those items.} \\
\texttt{REFUSE: \{'A', 'B', 'C', …\}} \\
\texttt{Explicitly reject opponent's proposal.} \\
\texttt{ARGUMENT: \{'...'\}} \\
\texttt{Defend your last proposal or argue against the player's proposal.} \\
\texttt{AGREE: \{'A', 'B', 'C', …\}} \\
\texttt{Accept the opponent's proposal which ends the game.} \\
\texttt{STRATEGIC REASONING: \{'...'\}} \\
\texttt{	Describe your strategic reasoning or anticipation explaining your choice of action. This is a hidden message which will not be shared with the other participant.} \\
\\ 
\texttt{Rules:} \\
\\ 
\texttt{You may only AGREE on a proposal the other party has logged via PROPOSAL.} \\
\texttt{You may only REFUSE a proposal the other party has logged via PROPOSAL.} \\
\texttt{Total effort of any PROPOSAL or AGREE set must be ≤ LIMIT.} \\
\texttt{Do NOT reveal your hidden importance scores.} \\
\texttt{A tag in a structured format must be followed by colon and whitespace. The argument must be a python set containing 0 or more strings.} \\
\texttt{So, it must be of the form TAG: \{...\}} \\
\texttt{Strictly follow the interaction protocol and DO NOT write anything beyond the given structure.} \\
\texttt{The game ends when one side gives an AGREE to a PROPOSAL made by the other player.} \\
\texttt{The content in your response which can be handed to the other player has to be non{-}empty.} \\
\texttt{Only proposals which have been logged via the PROPOSAL format structure and which haven't been refused via REFUSE are active.} \\
\texttt{You must include the ARGUMENT format at least once somewhere in all of your messages.} \\
\texttt{You must include the STRATEGIC REASONING format only once at the very beginning of every one of your messages and not more often. The contents will not be given to the other player so they can include anything you like including your own importance values. Here you should reason step by step to come up with you next move.} \\
            }
        }
    }
    & \\ \\

    \theutterance \stepcounter{utterance}  
    & \multicolumn{4}{p{0.6\linewidth}}{
        \cellcolor[rgb]{0.9,0.9,0.9}{
            \makecell[{{p{\linewidth}}}]{
                \texttt{\tiny{[P1$\rangle$GM]}}
                \texttt{<think>} \\
\texttt{Okay, let's see. I need to figure out the best way to approach this negotiation game. The goal is to maximize my own total importance while staying within the effort limit of 6938. The other player's values are unknown, so I have to make strategic decisions based on my own priorities.} \\
\\ 
\texttt{First, I should calculate the total importance for each item to see which ones are most valuable to me. Let me list out the items in order of their importance. The highest importance items are A08 (1096), C10 (1044), A25 (1005), A94 (972), B95 (1049), B34 (903), B96 (788), A87 (711), B64 (627), B15 (626), C79 (623), B19 (849), C72 (832), etc. Wait, actually, B95 has 1049, which is higher than C10's 1044. So maybe B95 is the top. Let me sort them properly.} \\
\\ 
\texttt{Wait, the importance values are: A08:1096, B95:1049, C10:1044, A25:1005, A94:972, B34:903, B96:788, A87:711, B64:627, B15:626, C79:623, C72:832, B19:849... Hmm, maybe I should sort them in descending order. Let me do that.} \\
\\ 
\texttt{Top items by importance:} \\
\texttt{1. A08: 1096} \\
\texttt{2. B95: 1049} \\
\texttt{3. C10: 1044} \\
\texttt{4. A25: 1005} \\
\texttt{5. A94: 972} \\
\texttt{6. B34: 903} \\
\texttt{7. B96: 788} \\
\texttt{8. C72: 832} \\
\texttt{9. B19: 849} \\
\texttt{10. A8} \\
            }
        }
    }
    & & \\ \\

    \theutterance \stepcounter{utterance}  
    & & \multicolumn{4}{p{0.6\linewidth}}{
        \cellcolor[rgb]{0.9,0.9,0.9}{
            \makecell[{{p{\linewidth}}}]{
                \texttt{\tiny{[P1$\langle$GM]}}
                \texttt{Your response did not start with the proper strategic reasoning tag at the very beginning of your response. The very first structured format must be of the form STRATEGIC REASONING: \{...\}. Try again.} \\
            }
        }
    }
    & \\ \\

    \theutterance \stepcounter{utterance}  
    & \multicolumn{4}{p{0.6\linewidth}}{
        \cellcolor[rgb]{0.9,0.9,0.9}{
            \makecell[{{p{\linewidth}}}]{
                \texttt{\tiny{[P1$\rangle$GM]}}
                \texttt{<think>} \\
\texttt{Okay, let me start over. The user wants me to make sure that my response begins with the STRATEGIC REASONING tag. In my previous attempt, I didn't include it, so I need to correct that.} \\
\\ 
\texttt{First, I need to outline my strategic reasoning. My goal is to maximize my total importance while staying within the effort limit. I should prioritize items with the highest importance per unit of effort. Let me calculate the importance per effort ratio for each item. For example, A08 has an importance of 1096 and effort 997, so 1096/997 ≈ 1.1. B95 is 1049/950 ≈ 1.1. C10 is 1044/945 ≈ 1.1. A25 is 1005/906 ≈ 1.1. A94 is 972/873 ≈ 1.1. These are all similar ratios. Maybe I should look for items with the highest importance first, even if their effort is high, as long as the total doesn't exceed the limit.} \\
\\ 
\texttt{I should start by selecting the top items by importance. Let's list them in order: A08, B95, C10, A25, A94, B34, B96, C72, B19, etc. Now, I need to add their efforts and see if they fit under 6938. Let's try adding them one by one.} \\
\\ 
\texttt{A08: 997} \\
\texttt{B95: 950 → total 1947} \\
\texttt{C10: 945 → 2892} \\
\texttt{A25: 906 → 3798} \\
\texttt{A94: 873 → 4671} \\
\texttt{B34: 804 → 5475} \\
\texttt{B96: 689 → 6164} \\
\texttt{C72: 733 → 6897} \\
\texttt{B19: 750 → 7647 (exceeds limit)} \\
\\ 
\texttt{So up to C72, total is 6897. Adding B19 would exceed. Next, check if there's} \\
            }
        }
    }
    & & \\ \\

    \theutterance \stepcounter{utterance}  
    & & \multicolumn{4}{p{0.6\linewidth}}{
        \cellcolor[rgb]{0.9,0.9,0.9}{
            \makecell[{{p{\linewidth}}}]{
                \texttt{\tiny{[P1$\langle$GM]}}
                \texttt{Your response did not start with the proper strategic reasoning tag at the very beginning of your response. The very first structured format must be of the form STRATEGIC REASONING: \{...\}. Try again.} \\
            }
        }
    }
    & \\ \\

    \theutterance \stepcounter{utterance}  
    & \multicolumn{4}{p{0.6\linewidth}}{
        \cellcolor[rgb]{0.9,0.9,0.9}{
            \makecell[{{p{\linewidth}}}]{
                \texttt{\tiny{[P1$\rangle$GM]}}
                \texttt{<think>} \\
\texttt{Okay, I need to make sure I start with the STRATEGIC REASONING tag. Let me outline my plan again.} \\
\\ 
\texttt{First, I'll calculate the importance per effort ratio for each item to prioritize those that give the most value for the effort. Then, I'll select items starting from the highest ratio, ensuring the total effort doesn't exceed 6938. Let's list the items by their importance per effort:} \\
\\ 
\texttt{1. A88: 109/10 = 10.9} \\
\texttt{2. A97: 149/50 = 2.98} \\
\texttt{3. B23: 103/4 = 25.75} \\
\texttt{4. B37: 156/57 ≈ 2.74} \\
\texttt{5. B43: 167/68 ≈ 2.46} \\
\texttt{6. C84: 169/70 ≈ 2.41} \\
\texttt{7. C51: 163/64 ≈ 2.55} \\
\texttt{8. C82: 155/56 ≈ 2.77} \\
\texttt{9. A28: 336/237 ≈ 1.42} \\
\texttt{10. C53: 232/133 ≈ 1.74} \\
\\ 
\texttt{Wait, maybe I should recalculate these ratios correctly. Let me check each item:} \\
\\ 
\texttt{{-} A88: 109/10 = 10.9} \\
\texttt{{-} B23: 103/4 = 25.75} \\
\texttt{{-} B37: 156/57 ≈ 2.74} \\
\texttt{{-} B43: 167/68 ≈ 2.46} \\
\texttt{{-} C84: 169/70 ≈ 2.41} \\
\texttt{{-} C51: 163/64 ≈ 2.55} \\
\texttt{{-} C82: 155/56 ≈ 2.77} \\
\texttt{{-} A28: 336/237 ≈ 1.42} \\
\texttt{{-} C53: 232/133 ≈ 1.74} \\
\texttt{{-} A97} \\
            }
        }
    }
    & & \\ \\

    \theutterance \stepcounter{utterance}  
    & & & \multicolumn{2}{p{0.3\linewidth}}{
        \cellcolor[rgb]{0.9,0.9,0.9}{
            \makecell[{{p{\linewidth}}}]{
                \texttt{\tiny{[GM$|$GM]}}
                \texttt{abort} \\
            }
        }
    }
    & & \\ \\

    \theutterance \stepcounter{utterance}  
    & & & \multicolumn{2}{p{0.3\linewidth}}{
        \cellcolor[rgb]{0.9,0.9,0.9}{
            \makecell[{{p{\linewidth}}}]{
                \texttt{\tiny{[GM$|$GM]}}
                \texttt{end game} \\
            }
        }
    }
    & & \\ \\

\end{supertabular}
}

\end{document}
